\section{QSPI Peripheral}

\subsection{Overview}

The QSPI peripheral provides a flexible SPI/QSPI master interface with support for:
\begin{itemize}
  \item Single, Dual and Quad SPI modes
  \item Command, Address, Dummy and Data phases
  \item Execute-In-Place (XIP) mode
  \item Configurable clock generation
  \item Interrupt-driven operation
  \item FIFO-based transmit and receive buffering
\end{itemize}

The design is split into:
\begin{itemize}
  \item QSPI Kernel (protocol engine, FSM, shifters)
  \item FIFO buffers (TX/RX)
  \item Register interface (Wishbone via BPI)
\end{itemize}

\subsection{Architecture}

The peripheral consists of three major blocks:

\begin{enumerate}
  \item \textbf{Register Interface (Top Level)}
  \item \textbf{QSPI Kernel}
  \item \textbf{FIFO Subsystem}
\end{enumerate}

The kernel is controlled via memory-mapped registers and interacts with the FIFOs for data transfer.

\subsection{Register Model}

All registers are 64-bit wide.

\begin{itemize}
  \item CTRL: control bits (start, flush, mode configuration)
  \item CMD: command opcode (8-bit)
  \item ADDR: address (32-bit used)
  \item LEN: transfer length in bytes
  \item DUMMY: dummy cycles
  \item CLKDIV: clock divider
  \item TIMEOUT: timeout counter
  \item STATUS: runtime status flags
  \item FIFOSTAT: FIFO levels and flags
  \item XIPMODE: XIP mode bits
  \item Interrupt registers: ISR, RIS, IMSC, MIS, ICR
\end{itemize}

\subsection{Finite State Machine (FSM)}

\subsubsection{States}

\begin{itemize}
  \item IDLE
  \item CMD
  \item ADDR
  \item DUM
  \item DAT
  \item DONE
\end{itemize}

\subsubsection{State Transitions}

\begin{itemize}
  \item IDLE → CMD (normal start)
  \item IDLE → ADDR (XIP active)
  \item CMD → ADDR
  \item ADDR → DUM or DAT
  \item DUM → DAT
  \item DAT → DONE
  \item DONE → IDLE
\end{itemize}

\subsubsection{Important Behavior}

\begin{itemize}
  \item CMD phase is always Single-SPI
  \item XIP mode skips CMD phase
  \item DONE is a \textbf{single-cycle pulse}
  \item DAT termination depends on TX/RX mode and bit counter
\end{itemize}

\subsection{Clock Generation}

The SPI clock is generated as:

\[
f_{SCK} = \frac{f_{clk}}{2 \cdot (CLKDIV + 1)}
\]

Drive/sample signals occur \textbf{one system clock before} the actual SCK edge:
\begin{itemize}
  \item Drive: before falling edge (Mode 0)
  \item Sample: before rising edge
\end{itemize}

\subsection{Data Path}

\subsubsection{Bus Width Selection}

The number of bits per SCK cycle:
\begin{itemize}
  \item 1-bit (single)
  \item 2-bit (dual)
  \item 4-bit (quad)
\end{itemize}

\subsubsection{Shift Registers}

\paragraph{TX}
\begin{itemize}
  \item Loaded before DAT phase
  \item Shifted MSB-first
  \item Shift width depends on bus mode
\end{itemize}

\paragraph{RX}
\begin{itemize}
  \item Sampled on SCK sample events
  \item Shifted into register
\end{itemize}

\subsubsection{Critical Timing Fix (RX Path)}

Due to non-blocking assignment timing:
\begin{itemize}
  \item RX data is captured one cycle after kernel write signal
  \item A two-stage pipeline ensures correct data transfer to FIFO
\end{itemize}

\subsection{FIFO Subsystem}

\subsubsection{Structure}

\begin{itemize}
  \item Synchronous single-clock FIFO
  \item Depth: parameterizable (default 16)
  \item Width: 64-bit
  \item Implemented using registers (not SRAM)
\end{itemize}

Internal elements:
\begin{itemize}
  \item Write pointer
  \item Read pointer
  \item Fill counter
  \item Memory array
\end{itemize}

\subsubsection{Function}

\paragraph{Write Operation}
\begin{itemize}
  \item Allowed if FIFO not full
  \item Also allowed when simultaneous read occurs
\end{itemize}

\paragraph{Read Operation}
\begin{itemize}
  \item Allowed if FIFO not empty
  \item Also allowed when simultaneous write occurs
\end{itemize}

\paragraph{Simultaneous Read/Write}
\begin{itemize}
  \item No change in fill level
  \item Both pointers advance
\end{itemize}

\paragraph{Fill Level Tracking}
\begin{itemize}
  \item Counter tracks entries (0..DEPTH)
  \item Used for interrupt generation
\end{itemize}

\paragraph{Status Signals}
\begin{itemize}
  \item full / empty
  \item half full / half empty
  \item almost full / almost empty
\end{itemize}

\paragraph{Read Latency}
\begin{itemize}
  \item Output is combinational from memory
  \item Effective latency: 1 cycle (due to pointer update)
\end{itemize}

\paragraph{Flush}
\begin{itemize}
  \item Synchronous reset of pointers and counter
\end{itemize}

\subsection{Interrupt System}

\subsubsection{Sources}

\begin{itemize}
  \item DONE
  \item RX HALF
  \item TX HALF
  \item ERROR
  \item TIMEOUT
\end{itemize}

\subsubsection{RIS Register}

\begin{itemize}
  \item Latching behavior
  \item Set by:
    \begin{itemize}
      \item hardware events (edge detection)
      \item software (ISR write)
    \end{itemize}
  \item Cleared by:
    \begin{itemize}
      \item ICR write
    \end{itemize}
\end{itemize}

\subsubsection{Important Property}

Interrupts are generated on \textbf{edges}, not levels.

\subsection{XIP Mode}

\begin{itemize}
  \item Activated after DONE when enabled
  \item Skips CMD phase
  \item Mode bits are driven during first dummy cycle
\end{itemize}

\subsection{Error Handling}

Errors occur when:
\begin{itemize}
  \item TX underflow
  \item RX overflow
\end{itemize}

\subsection{Timeout}

\begin{itemize}
  \item Counter decrements during active transfer
  \item Timeout flag set when reaching zero
\end{itemize}

\subsection{Status Signals}

\begin{itemize}
  \item busy
  \item done (pulse)
  \item error
  \item timeout
  \item xip\_active
\end{itemize}